\documentclass{if-beamer}
\usepackage[english]{babel}
\usepackage[absolute,overlay]{textpos}
\usepackage{booktabs}
\usepackage{tikz}
\usetikzlibrary{shapes.geometric, arrows.meta, positioning, calc, fit, backgrounds}

% --------------------------------------------------- %
%                  Presentation Info                  %
% --------------------------------------------------- %
\title[CS399P Mid-Evaluation]{\Large\textbf{Development of a Simulation and Benchmarking Platform for AI-Based Image-Guided Medical Robotics}}
\subtitle{\normalsize\textbf{Major Technical Project (MTP) --- Mid-Evaluation}}

\author[]{
    \textbf{Rishang Yadav} \quad (\textbf{Roll No: B23173})\\[0.25em]
    \small Supervisor: \textbf{Dr. Deepak Raina}\\[0.15em]
    \footnotesize School of Computing and Electrical Engineering (SCEE)
}

\institute[IIT Mandi]{
    \small Indian Institute of Technology (IIT) Mandi
}

\date{October 2026}
\subject{CS399P MTP Mid-Evaluation Presentation}

% Logo configuration: Fallback vector badge if iitlogo.png is absent
\logoleft{%
  \IfFileExists{iitlogo.png}{%
    \includegraphics[width=2.8cm]{iitlogo.png}%
  }{%
    \begin{tikzpicture}
      \draw[fill=iitmaroon!12, draw=iitmaroon, thick, rounded corners=3pt] (0,0) rectangle (2.8, 1.0);
      \node[align=center, font=\scriptsize\bfseries, text=iitmaroon] at (1.4, 0.5) {IIT MANDI\\[0.1em]CS399P};
    \end{tikzpicture}%
  }%
}

\begin{document}

% --------------------------------------------------- %
% Slide 1: Title Page
% --------------------------------------------------- %
\begin{frame}
  \titlepage
  \note{
    \textbf{Speaker Notes (Title Page):}\\
    - Introduce yourself: Rishang Yadav, Roll No. B23173, working under Dr. Deepak Raina in SCEE, IIT Mandi.\\
    - Project Title: Developing a high-performance simulation platform and autonomous ultrasound framework for next-generation medical robotics.\\
    - Set the stage: This mid-evaluation covers the development of our simulation bridge, physical acoustic pipeline, and the architectural foundation for AI training.
  }
\end{frame}

% --------------------------------------------------- %
% Slide 2: Motivation & Problem Statement
% --------------------------------------------------- %
\begin{frame}{Motivation: The Healthcare AI Dilemma}
    \begin{columns}[T]
        \begin{column}{0.48\textwidth}
            \begin{block}{Why Real-World Data Fails}
                \begin{itemize}
                    \item Modern AI requires \textbf{thousands of diverse trials} to master contact-rich interaction.
                    \item Collecting physical patient datasets is:
                    \begin{itemize}
                        \item \textbf{High Risk:} Cannot risk patient safety.
                        \item \textbf{Resource Intensive:} Clinical scans are costly and time-limited.
                        \item \textbf{Hard to Standardize:} Anatomical variation across patients.
                    \end{itemize}
                    \item \textbf{Physics-faithful simulation is the mandatory pre-condition for medical robot learning.}
                \end{itemize}
            \end{block}
        \end{column}
        \begin{column}{0.48\textwidth}
            \begin{alertblock}{Ultrasound: The Benchmark Task}
                \begin{itemize}
                    \item Ultrasound is non-invasive and radiation-free.
                    \item However, image acquisition is \textbf{strictly coupled to physical contact}:
                    \begin{itemize}
                        \item Sub-optimal angle $\to$ \textbf{acoustic shadow / blackout}.
                        \item Insufficient normal force $\to$ \textbf{acoustic decoupling}.
                        \item Excessive pressure $\to$ \textbf{tissue deformation / pain}.
                    \end{itemize}
                    \item Requires tight, simultaneous coordination of \textbf{probe pose, contact force, and image interpretation}.
                \end{itemize}
            \end{alertblock}
        \end{column}
    \end{columns}
    \note{
        \textbf{Speaker Notes (Motivation):}\\
        - Point out why standard robotics data collection pipelines fail in healthcare: trial-and-error on living patients is impossible.\\
        - Explain why robotic ultrasound is our chosen benchmark: unlike optical cameras where pulling back 2 cm still gives an image, ultrasound requires steady normal force; losing contact yields complete acoustic noise.\\
        - Therefore, simulation must accurately model physical contact dynamics and acoustic wave formation together.
    }
\end{frame}

% --------------------------------------------------- %
% Slide 3: Two Project Objectives (Visual Flow Diagram)
% --------------------------------------------------- %
\begin{frame}{Core Objectives}
    \centering
    \begin{tikzpicture}[
        node distance=0.3cm and 0.5cm,
        objbox1/.style={rectangle, draw=iitmaroon, fill=iitmaroon!8, thick, rounded corners=4pt, align=center, font=\small, inner sep=4pt, text width=6.0cm},
        objbox2/.style={rectangle, draw=techblue, fill=techblue!8, thick, rounded corners=4pt, align=center, font=\small, inner sep=4pt, text width=6.0cm},
        subbox/.style={rectangle, draw=darktext!30, fill=white, rounded corners=3pt, align=left, font=\scriptsize, inner sep=4pt, text width=5.6cm},
        flowarr/.style={-{Stealth[scale=1.1]}, very thick, draw=accentorange}
    ]
        % Left: Objective 1
        \node[objbox1] (obj1) {
            \textbf{\textcolor{iitmaroon}{OBJECTIVE 1: Build the Simulation Platform}}\\[0.15em]
            \footnotesize\textit{Realistic Virtual Environment for Medical Robotics}
        };
        
        \node[subbox, below=0.2cm of obj1] (sub1) {
            \textbullet\ \textbf{Tissue Contact \& Forces:} Realistic probe pressure at 1000 Hz\\
            \textbullet\ \textbf{Ultrasound Physics:} Real-time B-scans with shadows from CT\\
            \textbullet\ \textbf{Hospital Software Link:} Live streaming into 3D Slicer\\
            \textbullet\ \textbf{Safety Boundaries:} Virtual keep-out zones for sensitive organs
        };

        % Right: Objective 2
        \node[objbox2, right=1.6cm of obj1] (obj2) {
            \textbf{\textcolor{techblue}{OBJECTIVE 2: Train \& Benchmark AI Models}}\\[0.15em]
            \footnotesize\textit{Autonomous Scanning \& Evaluation Testbed}
        };
        
        \node[subbox, below=0.2cm of obj2] (sub2) {
            \textbullet\ \textbf{Active 3D Scanning:} Guide the probe to map full organs\\
            \textbullet\ \textbf{Smooth Control:} Learn safe probe trajectories (e.g., diffusion)\\
            \textbullet\ \textbf{Multimodal AI:} Test models using both images and force feedback\\
            \textbullet\ \textbf{Benchmarking:} Measure scan speed and quality against baselines
        };

        % Dynamic Flow Connector
        \draw[flowarr] (obj1.east) -- node[above, font=\scriptsize\bfseries, text=accentorange, align=center] {Safe Training Data} 
                                      node[below, font=\tiny, text=darktext] {(Ultrasound + Forces)} (obj2.west);
    \end{tikzpicture}

    \vspace{0.25cm}
    \begin{block}{\centering \textbf{How the Two Objectives Connect}}
        \scriptsize
        \textbf{Objective 1} builds the safe, physics-based simulator to generate realistic ultrasound and force data. \textbf{Objective 2} uses this simulated data to train, test, and benchmark autonomous scanning algorithms before deploying on physical robots.
    \end{block}
    \note{
        \textbf{Speaker Notes (Objectives):}\\
        - Keep it simple: Our project has two straightforward goals.\\
        - First, we build a realistic simulator (Objective 1) because training robots on real patients is impossible. It models probe contact forces and generates live ultrasound images with realistic wave shadows directly inside 3D Slicer.\\
        - Second, we use this simulator as a safe training ground (Objective 2) to teach robots how to scan autonomously and compare different AI approaches fairly.\\
        - The first two months focused on getting the simulator working and setting up the active scanning formulation.
    }
\end{frame}

% --------------------------------------------------- %
% Slide 4: Literature Survey (Part 1): Simulators
% --------------------------------------------------- %
\begin{frame}{Literature Survey (1/2): Robotic Ultrasound \& Simulators}
    \begin{columns}[T]
        \begin{column}{0.48\textwidth}
            \begin{block}{\textbf{SonoGym (Ao et al., 2025)}}
                \begin{itemize}\itemsep=3pt
                    \item \textbf{Overview:} Benchmark for robotic ultrasound using NVIDIA IsaacLab and clinical CT scans.
                    \item \textbf{Core Strengths:} Fast GPU-parallel RL environment (IsaacLab) for 2D probe tracking and drilling.
                    \item \textcolor{accentorange}{\textbf{Key Limitations:}}
                    \begin{itemize}\itemsep=2pt
                        \item \textbf{Bypasses Contact Dynamics:} Assumes an ideal controller keeps probe glued to skin; no continuous force-coupled acoustic mechanics.
                        \item \textbf{Static Anatomy:} Patient model is rigid; zero tissue deformation under probe pressure.
                        \item \textbf{No Clinical UI:} Headless/IsaacLab RL setup; completely disconnected from clinical surgical suites (no 3D Slicer).
                        \item \textbf{2D GAN Ultrasound:} Lacks physical acoustic ray tracing (cannot compute wave attenuation or off-axis bone shadows).
                    \end{itemize}
                \end{itemize}
            \end{block}
        \end{column}

        \begin{column}{0.48\textwidth}
            \begin{block}{\textbf{Slicer Bridges (AMBF, SlicerROS2)}}
                \begin{itemize}\itemsep=3pt
                    \item \textbf{Overview:} Frameworks bridging physics simulators to 3D Slicer for medical visualization.
                    \item \textbf{Core Strengths:} Demonstrates the clinical value of 3D Slicer as an open-source surgical UI.
                    \item \textcolor{accentorange}{\textbf{Key Limitations:}}
                    \begin{itemize}\itemsep=2pt
                        \item \textbf{Middleware Complexity:} Relies on heavy ROS/ROS2 communication pipelines.
                        \item \textbf{Serialization Latency:} Multi-process ROS serialization introduces latency and jitter during high-rate control.
                        \item \textbf{Setup Overhead:} Requires multi-node Linux daemons.
                    \end{itemize}
                \end{itemize}
            \end{block}
        \end{column}
    \end{columns}
    \note{
        \textbf{Speaker Notes (Literature Survey - Part 1):}\\
        - In this first literature review slide, we focus on medical ultrasound simulation environments.\\
        - SonoGym is our primary reference: built on NVIDIA IsaacLab, it demonstrated GPU-parallel RL for probe tracking. However, it has fundamental limitations: (1) it explicitly assumes an ideal controller keeps contact glued to the skin, bypassing contact force physics; (2) its patient model is completely rigid without tissue deformation; (3) it has zero integration with clinical software like 3D Slicer; and (4) its 2D GAN ultrasound lacks physical wave propagation.\\
        - SlicerROS2 proved that 3D Slicer is the right medical interface, but running ROS2 introduces significant IPC latency and setup bloat on clinical computers.\\
        - These gaps directly guided our work: direct OpenIGTLink sockets and GPU-based physical ray acoustics.
    }
\end{frame}

% --------------------------------------------------- %
% Slide 5: Literature Survey (Part 2): Surgical AI
% --------------------------------------------------- %
\begin{frame}{Literature Survey (2/2): Surgical Learning \& World Models}
    \begin{columns}[T]
        \begin{column}{0.48\textwidth}
            \begin{block}{\textbf{Surgical VLAs: SRT \& SRT-H (2024)}}
                \begin{itemize}\itemsep=3pt
                    \item \textbf{Overview:} Vision-Language-Action models and diffusion policies for surgical manipulation.
                    \item \textbf{Key Takeaways:}
                    \begin{itemize}\itemsep=2pt
                        \item Demonstrates that action chunking and language conditioning yield smooth, reactive motions.
                        \item High-level language commands enable intuitive supervisory control.
                    \end{itemize}
                    \item \textcolor{accentorange}{\textbf{Key Limitations:}}
                    \begin{itemize}\itemsep=2pt
                        \item Tested exclusively on rigid manipulation (peg-transfer and block tasks).
                        \item Lacks continuous physical contact compliance.
                    \end{itemize}
                \end{itemize}
            \end{block}
        \end{column}

        \begin{column}{0.48\textwidth}
            \begin{block}{\textbf{Surgical World Models (Cosmos-H, SAW)}}
                \begin{itemize}\itemsep=3pt
                    \item \textbf{Overview:} Video diffusion world models predicting future surgical scenes conditioned on robot actions.
                    \item \textbf{Key Takeaways:}
                    \begin{itemize}\itemsep=2pt
                        \item Demonstrates the ability of generative models to emulate procedural workflows.
                    \end{itemize}
                    \item \textcolor{accentorange}{\textbf{Key Limitations:}}
                    \begin{itemize}\itemsep=2pt
                        \item \textbf{Geometric Hallucination:} Generative video models hallucinate anatomical boundaries without rigid physical constraints.
                        \item \textbf{Computational Burden:} High inference latency.
                    \end{itemize}
                \end{itemize}
            \end{block}
        \end{column}
    \end{columns}
    \note{
        \textbf{Speaker Notes (Literature Survey - Part 2):}\\
        - In this second literature slide, we examine advanced surgical AI models.\\
        - Surgical Robot Transformers (SRT and SRT-H) show that diffusion action chunking works brilliantly for surgical robots, but their tasks were purely rigid (peg transfer). They do not incorporate contact-dependent imaging like ultrasound.\\
        - Surgical World Models like Cosmos-H show the promise of generative modeling, but pure video diffusion hallucinates boundaries and cannot enforce physical collision safety.\\
        - This motivates our hybrid direction: grounded physics and ray acoustics combined with generative action policies.
    }
\end{frame}

% --------------------------------------------------- %
% Slide 6: What We Are Doing Differently Than Baselines (Visual Comparison Diagram)
% --------------------------------------------------- %
\begin{frame}{What We Are Doing Differently Than Baselines}
    \centering
    \begin{tikzpicture}[
        node distance=0.2cm and 0.4cm,
        colhead/.style={rectangle, draw=darktext!50, fill=darktext!8, font=\small\bfseries, align=center, rounded corners=3pt, text width=4.3cm, minimum height=0.5cm},
        basebox/.style={rectangle, draw=iittwilight!40, fill=iittwilight!5, font=\scriptsize, align=left, rounded corners=3pt, text width=4.1cm, inner sep=4pt},
        ourbox/.style={rectangle, draw=techblue!80, fill=techblue!8, font=\scriptsize, align=left, rounded corners=3pt, text width=4.1cm, inner sep=4pt},
        parr/.style={-{Stealth[scale=0.9]}, thick, draw=accentorange}
    ]
        % Column 1: Clinical Interface
        \node[colhead] (c1) {\textbf{1. Clinical Interface}};
        \node[basebox, below=0.15cm of c1] (b1) {
            \textcolor{iittwilight}{\textbf{Established Baselines:}}\\
            \textbullet\ ROS / ROS2 middleware (SlicerROS2)\\
            \textbullet\ Multi-process IPC serialization\\
            \textbullet\ Headless / Omniverse GUI (SonoGym)
        };
        \node[ourbox, below=0.3cm of b1] (o1) {
            \textcolor{techblue}{\textbf{Our Proposed Approach:}}\\
            \textbullet\ \textbf{Direct OpenIGTLink TCP Sockets}\\
            \textbullet\ Hospital-standard surgical protocol\\
            \textbullet\ 1000 Hz physics $\leftrightarrow$ 60 FPS in-place VTK
        };
        \draw[parr] (b1.south) -- node[right, font=\tiny\bfseries, text=accentorange] {Direct Socket} (o1.north);

        % Column 2: Ultrasound Generation
        \node[colhead, right=0.4cm of c1] (c2) {\textbf{2. Ultrasound Generation}};
        \node[basebox, below=0.15cm of c2] (b2) {
            \textcolor{iittwilight}{\textbf{Established Baselines:}}\\
            \textbullet\ 2D pix2pix GAN synthesis (SonoGym)\\
            \textbullet\ Image-space texture prediction\\
            \textbullet\ Learned from specific ex-vivo scans
        };
        \node[ourbox, below=0.3cm of b2] (o2) {
            \textcolor{techblue}{\textbf{Our Proposed Approach:}}\\
            \textbullet\ \textbf{Model-Based GPU Ray Shader}\\
            \textbullet\ Real-time wave attenuation: $e^{-\alpha f z}$\\
            \textbullet\ Physical impedance reflections \& speckle
        };
        \draw[parr] (b2.south) -- node[right, font=\tiny\bfseries, text=accentorange] {Wave Physics} (o2.north);

        % Column 3: 3D Scanning & Planning
        \node[colhead, right=0.4cm of c2] (c3) {\textbf{3. 3D Scanning \& Planning}};
        \node[basebox, below=0.15cm of c3] (b3) {
            \textcolor{iittwilight}{\textbf{Established Baselines:}}\\
            \textbullet\ 2D image-plane visual tracking\\
            \textbullet\ Open-loop blind raster/zigzag sweeps\\
            \textbullet\ Susceptible to acoustic shadow voids
        };
        \node[ourbox, below=0.3cm of b3] (o3) {
            \textcolor{techblue}{\textbf{Our Proposed Approach:}}\\
            \textbullet\ \textbf{Active 3D Voxel Reconstruction}\\
            \textbullet\ Submodular information-gain planner\\
            \textbullet\ Formulated to steer around bone shadows
        };
        \draw[parr] (b3.south) -- node[right, font=\tiny\bfseries, text=accentorange] {Active NBV} (o3.north);
    \end{tikzpicture}

    \vspace{0.2cm}
    \note{
        \textbf{Speaker Notes (What We Do Differently - Visual Comparison):}\\
        - Walk the committee through the three vertical visual comparisons:\\
        - 1. Interface: Rather than dragging heavy ROS middleware or operating headlessly in IsaacLab, we stream directly into 3D Slicer over OpenIGTLink binary TCP sockets.\\
        - 2. Ultrasound: Rather than relying on a 2D GAN that hallucinates boundaries and lacks wave physics, our GPU shader computes real acoustic impedance, depth attenuation, and speckle noise directly from CT tensors.\\
        - 3. Planning: Rather than simple 2D tracking or blind raster sweeping, our pipeline accumulates 3D voxels with confidence weighting and optimizes Next-Best-Views to eliminate shadows.
    }
\end{frame}

% --------------------------------------------------- %
% Slide 7: Work Accomplished: Slicer Surgical Bridge
% --------------------------------------------------- %
\begin{frame}{Work Accomplished: Slicer-PyBullet Surgical Bridge}
    \centering
    \begin{tikzpicture}[
        node distance=0.6cm and 0.8cm,
        slicerbox/.style={rectangle, draw=techblue, fill=techblue!10, thick, rounded corners=3pt, align=center, font=\scriptsize, inner sep=4pt, text width=3.3cm},
        pybox/.style={rectangle, draw=iitmaroon, fill=iitmaroon!10, thick, rounded corners=3pt, align=center, font=\scriptsize, inner sep=4pt, text width=3.3cm},
        netbox/.style={rectangle, draw=accentorange, fill=accentorange!12, dashed, thick, rounded corners=3pt, align=center, font=\scriptsize\bfseries, inner sep=4pt, text width=3.4cm},
        netarr/.style={-{Stealth[scale=0.9]}, thick, draw=accentorange}
    ]
        % Slicer Side
        \node[slicerbox] (slicer) {
            \textbf{3D Slicer Interface}\\
            \textit{\tiny (Surgeon's Clinical Viewport)}\\[0.2em]
            \textbullet\ \textbf{3D Viewport:} Shows robot \& anatomy\\
            \textbullet\ \textbf{In-Place VTK:} Lag-free 60 FPS visual\\
            \textbullet\ \textbf{Force Vectors:} Visual pressure arrows\\
            \textbullet\ \textbf{Safety Zones:} Surgeon keep-out limits
        };

        % OpenIGTLink Bridge
        \node[netbox, right=0.8cm of slicer] (bridge) {
            \textbf{OpenIGTLink Protocol}\\
            \textit{\tiny (Hospital-Standard Network Link)}\\[0.3em]
            \textcolor{techblue}{$\leftarrow$ User Targets \& Safety Zones}\\
            \textcolor{techblue}{\tiny \textit{(Sent from clinician UI to simulator)}}\\[0.2em]
            \textcolor{iitmaroon}{$\rightarrow$ Live Robot Poses \& Contact Forces}\\
            \textcolor{iitmaroon}{\tiny \textit{(Streamed from simulator to UI at 60 FPS)}}
        };

        % PyBullet Side
        \node[pybox, right=0.8cm of bridge] (pybullet) {
            \textbf{Physics Server (PyBullet)}\\
            \textit{\tiny (Virtual Robot \& Patient World)}\\[0.2em]
            \textbullet\ \textbf{1000 Hz Loop:} Stable contact dynamics\\
            \textbullet\ \textbf{Soft Tissue:} Skin deform under probe\\
            \textbullet\ \textbf{Force Sensing:} Computes contact pressure\\
            \textbullet\ \textbf{Virtual Fixtures:} Physical barrier stops
        };

        \draw[netarr] (slicer.east) -- (bridge.west);
        \draw[netarr] (bridge.east) -- (pybullet.west);
    \end{tikzpicture}

    \vspace{0.25cm}
    \begin{columns}[T]
        \begin{column}{0.32\textwidth}
            \begin{block}{\centering \textbf{Decoupled Multi-Rate Timing}}
                \footnotesize
                \textbullet\ \textbf{Fast Physics (1000 Hz):} Prevents contact explosion.\\
                \textbullet\ \textbf{Smooth Display (60 FPS):} Clinical UI responsiveness.\\
                \textbullet\ Rendering hiccups never crash physical simulation.
            \end{block}
        \end{column}
        \begin{column}{0.32\textwidth}
            \begin{block}{\centering \textbf{In-Place Memory Updates}}
                \footnotesize
                \textbullet\ Overwrites existing graphics memory buffers directly.\\
                \textbullet\ Eliminates deleting and recreating 3D objects.\\
                \textbullet\ Guarantees zero GUI stutter during long scans.
            \end{block}
        \end{column}
        \begin{column}{0.32\textwidth}
            \begin{block}{\centering \textbf{Virtual Safety Fixtures}}
                \footnotesize
                \textbullet\ Surgeon delineates risk zones directly on 3D CT.\\
                \textbullet\ Exported as hard physical obstacles in PyBullet.\\
                \textbullet\ Robot is physically blocked from sensitive anatomy.
            \end{block}
        \end{column}
    \end{columns}
    \note{
        \textbf{Speaker Notes (Slicer Bridge Architecture):}\\
        - The flow diagram illustrates the architecture: PyBullet calculates contact dynamics and kinematics, transmitting binary OpenIGTLink messages over TCP directly to 3D Slicer.\\
        - Multi-rate decoupling: Contact mechanics require high frequency to prevent numerical instability, whereas visualization requires 60 FPS. Decoupling them ensures rendering hiccups never crash physical simulation.\\
        - In-place memory: Rather than destroying and re-creating VTK scene nodes every frame, we mutate the coordinate buffer in place, preventing memory churn.\\
        - Virtual fixtures: Allows clinicians to delineate sensitive anatomical structures (e.g., spinal cord) which PyBullet enforces as physical keep-out obstacles.
    }
\end{frame}

% --------------------------------------------------- %
% Slide 8: Work Accomplished: 3D Ultrasound Pipeline
% --------------------------------------------------- %
\begin{frame}{Work Accomplished: 3D Ultrasound Reconstruction Pipeline}
    \centering
    \begin{tikzpicture}[
        node distance=0.4cm and 0.5cm,
        box/.style={rectangle, draw=techblue, fill=techblue!8, thick, rounded corners=3pt, align=center, font=\scriptsize, inner sep=3pt, text width=2.4cm, minimum height=1.0cm},
        rbox/.style={rectangle, draw=iitmaroon, fill=iitmaroon!8, thick, rounded corners=3pt, align=center, font=\scriptsize, inner sep=3pt, text width=2.4cm, minimum height=1.0cm},
        pbox/.style={rectangle, draw=accentorange, fill=accentorange!10, thick, rounded corners=3pt, align=center, font=\scriptsize, inner sep=3pt, text width=2.4cm, minimum height=1.0cm},
        kbox/.style={rectangle, draw=iittwilight, fill=iittwilight!8, thick, rounded corners=3pt, align=center, font=\scriptsize, inner sep=3pt, text width=2.4cm, minimum height=1.0cm},
        arr/.style={-{Stealth[scale=0.85]}, thick, draw=iittwilight}
    ]
        % Top Row
        \node[box] (ct) {\textbf{Patient CT Scan}\\[0.15em]\tiny 3D anatomical volume data};
        \node[box, right=of ct] (gpu) {\textbf{GPU Reslicer}\\[0.15em]\tiny Cuts 2D slice at probe pose};
        \node[box, right=of gpu] (ray) {\textbf{Acoustic Shader}\\[0.15em]\tiny Simulates sound decay \& noise};
        \node[rbox, right=of ray] (seg) {\textbf{Confidence Map}\\[0.15em]\tiny Flags bone shadows \& blind spots};

        % Bottom Row
        \node[kbox, below=0.5cm of ct] (ik) {\textbf{Robot Controller}\\[0.15em]\tiny Moves 7-DoF arm to target};
        \node[pbox, right=of ik] (nbv) {\textbf{Active View Planner}\\[0.15em]\tiny Picks best angle to see bone};
        \node[pbox, right=of nbv] (pnet) {\textbf{Uncertainty Map}\\[0.15em]\tiny Spots unmapped / dark areas};
        \node[rbox, right=of pnet] (accum) {\textbf{3D Voxel Fusion}\\[0.15em]\tiny Stitches 2D slices into 3D volume};

        % Connections
        \draw[arr] (ct) -- (gpu);
        \draw[arr] (gpu) -- (ray);
        \draw[arr] (ray) -- (seg);
        \draw[arr] (seg.south) -- ++(0,-0.25) -| (accum.north);
        \draw[arr] (accum) -- (pnet);
        \draw[arr] (pnet) -- (nbv);
        \draw[arr] (nbv) -- (ik);
        \draw[arr, dashed, draw=accentorange, thick] (ik.north) -- (ct.south);
    \end{tikzpicture}

    \vspace{0.2cm}
    \begin{columns}[T]
        \begin{column}{0.32\textwidth}
            \begin{block}{\centering \textbf{GPU Sound Wave Physics}}
                \footnotesize
                \textbullet\ \textbf{Sub-2ms Slicing:} PyTorch trilinear sampling.\\
                \textbullet\ \textbf{Exponential Decay:} Sound dims with tissue depth.\\
                \textbullet\ \textbf{Rayleigh Noise:} Authentic ultrasound speckle.\\
                \textbullet\ Physics-based alternative to black-box GANs.
            \end{block}
        \end{column}
        \begin{column}{0.32\textwidth}
            \begin{block}{\centering \textbf{3D Volumetric Fusion}}
                \footnotesize
                \textbullet\ Stitches 2D frames into a global 3D voxel grid.\\
                \textbullet\ Uses tracked robot poses to place voxels.\\
                \textbullet\ \textbf{Confidence Weighing:} Blends clear pixels, discards shadows.\\
                \textbullet\ Builds a continuous 3D organ volume.
            \end{block}
        \end{column}
        \begin{column}{0.32\textwidth}
            \begin{block}{\centering \textbf{Active Shadow Avoidance}}
                \footnotesize
                \textbullet\ Prioritizes dark, unmapped regions.\\
                \textbullet\ Intelligently tilts probe to peer around bone.\\
                \textbullet\ Avoids wasting time re-scanning flat areas.\\
                \textbullet\ Formulated for complete organ coverage.
            \end{block}
        \end{column}
    \end{columns}
    \note{
        \textbf{Speaker Notes (Ultrasound Pipeline):}\\
        - Explain the closed loop in the flowchart: We extract 2D CT slices on the GPU in tensor memory, apply physical acoustic attenuation and speckle noise, and calculate confidence.\\
        - The 3D accumulator projects these slices into a global voxel space using tracked robot forward kinematics, building a continuous volume.\\
        - The active view planner evaluates spatial uncertainty to determine the next probe trajectory.\\
        - Note: This is an architectural formulation designed to resolve acoustic shadows; empirical verification is part of our future testing plan.
    }
\end{frame}

% --------------------------------------------------- %
% Slide 9: Autonomous Scanning Formulation
% --------------------------------------------------- %
\begin{frame}{Autonomous Scanning: Formulating Active View Selection}
    \begin{columns}[T]
        \begin{column}{0.48\textwidth}
            \begin{alertblock}{Limitations of Blind Raster Sweeps}
                \begin{itemize}
                    \item Baselines typically rely on \textbf{open-loop zigzag/raster sweeps}:
                    \begin{itemize}
                        \item Blindly traverses a planar bounding box.
                        \item Misses out-of-plane features, steep bone facets, and occluded contours.
                        \item Wastes scan time re-sweeping already reconstructed regions.
                    \end{itemize}
                    \item Real human anatomy is 3D, non-planar, and acoustically complex.
                \end{itemize}
            \end{alertblock}
        \end{column}
        \begin{column}{0.48\textwidth}
            \begin{block}{Active Exploration Formulation}
                \begin{itemize}
                    \item \textbf{Submodular Objective Design:}
                    \begin{itemize}
                        \item Re-scanning an already covered voxel yields diminishing information gain.
                        \item Formulated to balance new surface coverage against robot motion penalty.
                    \end{itemize}
                    \item \textbf{Expected Benefits:}
                    \begin{itemize}
                        \item Designed to steer probe tilt around acoustic obstacles (e.g., ribs).
                        \item Expected to reduce scan duration and eliminate shadow voids.
                    \end{itemize}
                    \item \textbf{Status:} Formulation implemented in simulation; systematic testing against baselines planned in future work.
                \end{itemize}
            \end{block}
        \end{column}
    \end{columns}
    \note{
        \textbf{Speaker Notes (Active Planning Formulation):}\\
        - Contrast blind raster scanning with active view planning: raster sweeps repeatedly scan flat regions while missing complex angled surfaces behind acoustic barriers.\\
        - Explain the principle of diminishing returns: re-scanning covered tissue should yield zero reward, encouraging the robot to explore unmapped areas.\\
        - Crucially clarify to the committee: This formulation is designed and implemented, but its comparative performance against baselines will be formally tested in the next project phase.
    }
\end{frame}

% --------------------------------------------------- %
% Slide 10: Technical Understanding: Three Core Decisions (Visual Decision Diagram)
% --------------------------------------------------- %
\begin{frame}{Technical Understanding}
    \centering
    \begin{tikzpicture}[
        node distance=0.2cm and 0.4cm,
        dcard/.style={rectangle, draw=darktext!40, fill=white, rounded corners=4pt, inner sep=4pt, text width=4.3cm, align=center},
        altbox/.style={rectangle, draw=iittwilight!40, fill=iittwilight!5, rounded corners=2pt, font=\scriptsize, inner sep=4pt, text width=4.0cm, align=left},
        chobox1/.style={rectangle, draw=iitmaroon!80!black, fill=iitmaroon!8, rounded corners=2pt, font=\scriptsize, inner sep=4pt, text width=4.0cm, align=left},
        chobox2/.style={rectangle, draw=techblue!80!black, fill=techblue!8, rounded corners=2pt, font=\scriptsize, inner sep=4pt, text width=4.0cm, align=left},
        chobox3/.style={rectangle, draw=accentorange!80!black, fill=accentorange!8, rounded corners=2pt, font=\scriptsize, inner sep=4pt, text width=4.0cm, align=left},
        dhead/.style={font=\small\bfseries, align=center, text width=4.1cm, minimum height=0.5cm}
    ]
        % Decision 1
        \node[dcard] (card1) {
            {\small\bfseries\textcolor{iitmaroon}{\textbf{1. Network Transport}}}\\[0.2em]
            \begin{tikzpicture}[node distance=0.15cm]
                \node[altbox] (r1) {
                    \textbf{\textcolor{iittwilight}{Alternative: ROS2 / DDS}}\\
                    Standard robotics middleware; introduces multi-node daemons and IPC serialization.
                };
                \node[chobox1, below=0.15cm of r1] (c1) {
                    \textbf{\textcolor{iitmaroon}{Design Choice: OpenIGTLink TCP}}\\
                    Native 3D Slicer protocol; direct, lightweight binary streaming between physics and UI.
                };
            \end{tikzpicture}
        };

        % Decision 2
        \node[dcard, right=0.4cm of card1] (card2) {
            {\small\bfseries\textcolor{techblue}{\textbf{2. Ultrasound Generation}}}\\[0.2em]
            \begin{tikzpicture}[node distance=0.15cm]
                \node[altbox] (r2) {
                    \textbf{\textcolor{iittwilight}{Alternative: 2D Generative GAN}}\\
                    Data-driven texture synthesis, but requires retraining per patient and lacks ray wave physics.
                };
                \node[chobox2, below=0.15cm of r2] (c2) {
                    \textbf{\textcolor{techblue}{Design Choice: GPU Ray Physics}}\\
                    Deterministic wave attenuation and acoustic impedance; evaluates zero-shot on any patient CT.
                };
            \end{tikzpicture}
        };

        % Decision 3
        \node[dcard, right=0.4cm of card2] (card3) {
            {\small\bfseries\textcolor{accentorange}{\textbf{3. View Planning Strategy}}}\\[0.2em]
            \begin{tikzpicture}[node distance=0.15cm]
                \node[altbox] (r3) {
                    \textbf{\textcolor{iittwilight}{Alternative: Heuristic Raster Sweeps}}\\
                    Standard scanning pattern; simple open-loop coverage, but susceptible to persistent shadows.
                };
                \node[chobox3, below=0.15cm of r3] (c3) {
                    \textbf{\textcolor{accentorange}{Design Choice: Submodular NBV}}\\
                    Evaluates marginal information gain across 3D voxels; formulated to steer around bone.
                };
            \end{tikzpicture}
        };
    \end{tikzpicture}

    \vspace{0.25cm}
    \note{
        \textbf{Speaker Notes (Technical Understanding - Decision Architecture):}\\
        - Walk the committee through our three core design choices:\\
        - Decision 1 (OpenIGTLink): Why not ROS2? 3D Slicer natively supports OpenIGTLink. It allows direct TCP binary socket streaming with sub-millisecond round-trip time, avoiding ROS serialization overhead on clinical workstations.\\
        - Decision 2 (Physics over GANs): GANs produce plausible textures but hallucinate boundaries and cannot physically model sound attenuation. Our GPU shader models real acoustic physics directly from CT densities.\\
        - Decision 3 (Submodular Planning): Volumetric coverage is non-Markovian; submodular information gain naturally captures diminishing returns.
    }
\end{frame}

% --------------------------------------------------- %
% Slide 11: Future Direction (Part 1): Testing & Control
% --------------------------------------------------- %
\begin{frame}{Future Direction (1/2): Testing \& Interactive Control}
    \begin{columns}[T]
        \begin{column}{0.48\textwidth}
            \begin{block}{1. Simulation Pipeline Testing \& Verification}
                \begin{itemize}\itemsep=3pt
                    \item \textbf{Multi-Rate Timing Stability:}
                    \begin{itemize}\itemsep=1pt
                        \item Stress-test the physics loop under rapid probe contact force transitions and varied penetration depths.
                        \item Verify collision stability across varied tissue stiffness parameters.
                    \end{itemize}
                    \item \textbf{GPU Reslicing Profiling:}
                    \begin{itemize}\itemsep=1pt
                        \item Profile slicing and acoustic ray shader throughput across diverse patient CT volume sizes.
                    \end{itemize}
                    \item \textbf{OpenIGTLink Streaming Integrity:}
                    \begin{itemize}\itemsep=1pt
                        \item Verify continuous packet transmission during prolonged sessions.
                    \end{itemize}
                \end{itemize}
            \end{block}
        \end{column}

        \begin{column}{0.48\textwidth}
            \begin{block}{2. Interactive Trajectory Initialization}
                \begin{itemize}\itemsep=3pt
                    \item \textbf{Clinician-in-the-Loop Setup:}
                    \begin{itemize}\itemsep=1pt
                        \item Connect 3D Slicer interactive markups (fiducial points and bounding regions) directly to the simulation.
                        \item Clinicians can delineate a region of interest (e.g., lumbar vertebra) directly on the 3D viewport.
                    \end{itemize}
                    \item \textbf{Kinematic Goal Resolution:}
                    \begin{itemize}\itemsep=1pt
                        \item PyBullet inverse kinematics (IK) automatically resolves clinician markups into initial collision-free robot poses.
                    \end{itemize}
                \end{itemize}
            \end{block}
        \end{column}
    \end{columns}
    \note{
        \textbf{Speaker Notes (Future Direction - Part 1):}\\
        - In the first half of our future roadmap, we focus on engineering verification and interactive clinical controls.\\
        - Testing is our top priority: we need to systematically verify contact stability under high forces, profile GPU slicing speed across varied CT resolutions, and confirm socket packet integrity.\\
        - In addition, we will integrate 3D Slicer interactive markups so a clinician can click or drag target points on patient anatomy, allowing PyBullet IK to initialize scanning trajectories automatically.
    }
\end{frame}

% --------------------------------------------------- %
% Slide 12: Future Direction (Part 2): Policy Training
% --------------------------------------------------- %
\begin{frame}{Future Direction (2/2): Policy Training \& AI Benchmarking}
    \begin{columns}[T]
        \begin{column}{0.48\textwidth}
            \begin{block}{3. AI Policy Training in Simulation}
                \begin{itemize}\itemsep=3pt
                    \item \textbf{Generative Action Chunking:}
                    \begin{itemize}\itemsep=1pt
                        \item Train Continuous Flow Matching (CFM) and diffusion policies for smooth 6-DoF probe velocity generation.
                    \end{itemize}
                    \item \textbf{Active Exploration Optimization:}
                    \begin{itemize}\itemsep=1pt
                        \item Train policies using submodular information-gain rewards to dynamically avoid acoustic shadows.
                    \end{itemize}
                    \item \textbf{Multi-Patient Training:}
                    \begin{itemize}\itemsep=1pt
                        \item Train across multiple patient abdominal/spinal CT volumes to promote generalization.
                    \end{itemize}
                \end{itemize}
            \end{block}
        \end{column}

        \begin{column}{0.48\textwidth}
            \begin{block}{4. Comparative Benchmarking Suite}
                \begin{itemize}\itemsep=3pt
                    \item \textbf{Systematic Baseline Comparison:}
                    \begin{itemize}\itemsep=1pt
                        \item Empirically evaluate active planning against standard raster/zigzag sweeps in simulation.
                    \end{itemize}
                    \item \textbf{Evaluation Metrics (Expected):}
                    \begin{itemize}\itemsep=1pt
                        \item Reconstructed volume completeness percentage.
                        \item Acoustic shadow void reduction.
                        \item Total trajectory length and scan duration.
                    \end{itemize}
                    \item \textbf{VLA Observation Testbed:}
                    \begin{itemize}\itemsep=1pt
                        \item Structure observation spaces (B-mode ultrasound, contact force history, text prompts) for surgical VLA benchmarking.
                    \end{itemize}
                \end{itemize}
            \end{block}
        \end{column}
    \end{columns}
    \note{
        \textbf{Speaker Notes (Future Direction - Part 2):}\\
        - In the second half of our future roadmap, we focus on policy learning and establishing our simulation benchmarking suite.\\
        - We will implement diffusion and flow matching policies for smooth action chunking, and train them across diverse patient CT scans.\\
        - We will then run systematic comparative benchmarks: measuring whether our active view planning successfully reduces acoustic shadows and scan time compared to standard raster sweeps.\\
        - Finally, the platform will serve as an open benchmark for evaluating multimodal Vision-Language-Action models.
    }
\end{frame}

% --------------------------------------------------- %
% Slide 13: Conclusion & Committee Discussion
% --------------------------------------------------- %
\begin{frame}[c]{Conclusion \& Discussion}
    \centering
    \begin{block}{Summary of Work Accomplished}
        Over the first two months, the essential simulation infrastructure—the zero-ROS Slicer bridge and the physically grounded 3D ultrasound pipeline was built.
    \end{block}

    \vspace{0.4cm}
    {\Huge \textbf{\textcolor{iitmaroon}{Thank You!}}}\\[0.6em]
    {\large \textbf{Questions \& Discussion}}\\[1.0em]
    \note{
        \textbf{Speaker Notes (Conclusion):}\\
        - Summarize the core message: Over the first two months, we developed the essential simulation infrastructure—the zero-ROS Slicer bridge and the physically grounded 3D ultrasound pipeline.\\
        - Reassure the committee that the platform is now ready for systematic pipeline testing and AI policy benchmarking in the subsequent phase.\\
        - Open the floor for questions, ready to defend architectural choices.
    }
\end{frame}

\end{document}
